\documentclass{article}
\usepackage{graphicx}
\usepackage{subcaption}
\usepackage{amsmath}
\usepackage{amsfonts}
\usepackage{amssymb}
\usepackage{hyperref}
\usepackage{geometry}

\usepackage{booktabs}
\usepackage{multirow}

\geometry{margin=1in}

\setlength{\parindent}{0pt}

\title{Using Physics Informed Classical Machine Learning for the Detection of Exoplanets}
\author{Esraaj Sarkar Gupta, Suyash Prakash, Anirudh Puri}

\begin{document}
\maketitle

\section{Introduction}


\section{Methodology}
We use the Kepler Object of Interest (KOI) dataset to train our model.

\subsection{Transit Depth Consistency}
Since the telescope makes 2D measurements, we assert that the stellar flux must be directly related to the area of the star that is observed. During a transit, a portion of the area of the star is blocked by the planet, leading to a decrease in the observed flux. The transit depth, which is the fractional decrease in flux during a transit, can be expressed as:
\begin{equation}
    \delta_\text{theoretical} = \frac{\text{Area}_\text{Planet}}{\text{Area}_\text{Star}} = \frac{\pi R_p^2}{\pi R_*^2}
\end{equation}

where $R_p$ is the radius of the planet and $R_*$ is the radius of the star. We can compare this theoretical transit depth with the observed transit depth from the light curve data to ensure consistency.

\begin{equation}
    \Delta\delta = \Bigg|\delta_{obs} - \delta_\text{theoretical}\Bigg|
\end{equation}

\subsection{Duration Consistency}


\subsection{Impact Parameter}

\section{Results}

\begin{table}[h]
\centering
\caption{10-Fold Stratified Cross-Validation Results by Dataset and Feature Set. Note: KOI and TESS utilized complete cases, while K2 utilized KNN-imputed data.}
\label{tab:cv_results_full}
\begin{tabular}{llcc}
\toprule
\textbf{Dataset} & \textbf{Feature Set} & \textbf{Mean Accuracy} & \textbf{Mean PR-AUC} \\
\midrule
\multirow{3}{*}{KOI (k=10)}
& Ablation Features & $0.9467 \pm 0.0074$ & $0.9905 \pm 0.0019$ \\
& Features 3 & $0.9550 \pm 0.0066$ & $0.9913 \pm 0.0027$ \\
& Features 4 & $0.9544 \pm 0.0064$ & $0.9918 \pm 0.0023$ \\
\midrule
\multirow{3}{*}{TESS (k=10)}
& Ablation Features & $0.8481 \pm 0.0082$ & $0.5729 \pm 0.0405$ \\
& Features 3 & $0.8477 \pm 0.0107$ & $0.5774 \pm 0.0491$ \\
& Features 4 & $0.8478 \pm 0.0126$ & $0.5933 \pm 0.0451$ \\
\midrule
\multirow{3}{*}{K2 Imputed (k=10)}
& Ablation Features & $0.7948 \pm 0.0239$ & $0.9022 \pm 0.0149$ \\
& Features 3 & $0.7948 \pm 0.0275$ & $0.8957 \pm 0.0118$ \\
& Features 4 & $0.7879 \pm 0.0290$ & $0.8855 \pm 0.0148$ \\
\bottomrule
\end{tabular}
\end{table}


\begin{table}[h]
\centering
\caption{K2 Dataset Performance: Native NaN Handling vs. KNN Imputation (10-Fold CV)}
\label{tab:k2_missing_data}
\begin{tabular}{llcc}
\toprule
\textbf{Model \& Data Strategy} & \textbf{Feature Set} & \textbf{Mean Accuracy} & \textbf{Mean PR-AUC} \\
\midrule
\multirow{3}{*}{HistGB (Native NaN)}
& Ablation Features & $0.8380 \pm 0.0174$ & $0.9375 \pm 0.0096$ \\
& Features 3 & $0.8246 \pm 0.0161$ & $0.9297 \pm 0.0104$ \\
& Features 4 & $0.8102 \pm 0.0181$ & $0.9046 \pm 0.0184$ \\
\midrule
\multirow{3}{*}{RandomForest (Imputed)}
& Ablation Features & $0.7963 \pm 0.0241$ & $0.9022 \pm 0.0167$ \\
& Features 3 & $0.7971 \pm 0.0291$ & $0.8962 \pm 0.0168$ \\
& Features 4 & $0.7946 \pm 0.0277$ & $0.8870 \pm 0.0197$ \\
\bottomrule
\end{tabular}
\end{table}

\begin{table}[h]
\centering
\caption{TESS Dataset Performance: Native NaN Handling vs. KNN Imputation (10-Fold CV)}
\label{tab:tess_missing_data}
\begin{tabular}{llcc}
\toprule
\textbf{Model \& Data Strategy} & \textbf{Feature Set} & \textbf{Mean Accuracy} & \textbf{Mean PR-AUC} \\
\midrule
\multirow{3}{*}{HistGB (Native NaN)}
& Ablation Features & $0.8575 \pm 0.0126$ & $0.5543 \pm 0.0493$ \\
& Features 3 & $0.8602 \pm 0.0108$ & $0.5569 \pm 0.0496$ \\
& Features 4 & $0.8609 \pm 0.0083$ & $0.5591 \pm 0.0404$ \\
\midrule
\multirow{3}{*}{RandomForest (Imputed)}
& Ablation Features & $0.8590 \pm 0.0113$ & $0.5445 \pm 0.0471$ \\
& Features 3 & $0.8602 \pm 0.0091$ & $0.5431 \pm 0.0457$ \\
& Features 4 & $0.8610 \pm 0.0076$ & $0.5608 \pm 0.0430$ \\
\bottomrule
\end{tabular}
\end{table}

% Might want to discuss the limitations of imputation in





\end{document}